\scanpage{345}
% @ top centre
\mspage{69}

% @ top through parenthetical NB
Que $\unsure{\mathbb W_0|W_0}$ satisfasse à cet ensemble d'axiomes se vérifie \deleted{de façon} par \unsure{AQT|A.Q.T.}. (\uline{NB} Qu'il satisfasse à $\mathrm{Loc}(1\ \mathrm{ter},2,3\mathrm{bis})$ est déjà conséquence du fait que $\mathbb W_0$ est de la forme $\deleted{\unsure{\mathbb D}}W_{\mathbb D}$, où $\mathbb D$ ($=\mathbb D_{\mathbb W_0}$ en l'occurrence) est un prédérivateur satisfaisant $\mathrm{Der}(3,4,5)$, ce qui est bien le cas pour $\mathbb D_{\mathbb W_0}$, qui est un vrai \uline{dérivateur}.)

% @ centre and lower half, factorization
Il faut vérifier que tout $f$ dans $\mathrm{Cat}$ se factorise en $f=pi$, avec $i\in\mathbb W_0$, $p$ une $\mathbb W_0$-fibration. Or si $f:X\longrightarrow Y$, $I=\pi_0(X)$, on prend $\overline X=\coprod_{\unsure{d|\alpha}\in I}Y_{f(d)}$, \unsure{d'où} une factorisation $f=pi$, où $\deleted{\unsure{i\in}}\pi_0(i)$ iso i.e. $i\in\mathbb W_0$, et où $p:\overline X\longrightarrow\deleted{\unsure{X|Y}}$ est, au dessus de la composante connexe $Y_j$ de $Y$, donné par le produit $\overline X\mid Y_j\simeq Y_j\times\underbrace{\pi_0(f)^{-1}(j)}_{\text{cat. discrète}}$, \inserted{$p$} \deleted{\unsure{$q$}} est bien une $\mathbb W_0$-fibration, car grâce à $\mathrm{Loc}(4)$ il suffit de le vérifier sur chaque $Y_j$, et là c'est déduit d'une flèche $I_j\unsure{\longleftrightarrow|\longrightarrow}e$ (une $\mathbb W_0$-fibration) par chgt de base $Y_j\longrightarrow e$.

% @ foot, beginning of next argument
\deleted{\unsure{on|un}}\unsure{/|,} D'autre part, que $\mathbb W_0$ ne \unsure{satisfait|satisfasse} \uline{pas} à $D_{\unsure{\mathbb W_\infty|W_\infty}}$ \deleted{\unsure{est}} a déjà été vu ($\overline{\mathrm I}$, p.87),
% archive sheet number: 344
